\documentclass[a4paper,11pt]{article}

% --- Pacchetti Fondamentali ---
\usepackage[utf8]{inputenc}
\usepackage[T1]{fontenc}
\usepackage[italian]{babel}
\usepackage{geometry}
\geometry{top=2.5cm, bottom=2.5cm, left=2.5cm, right=2.5cm}

% --- Font Moderni ---
\usepackage{lmodern}
\usepackage{helvet}
\renewcommand{\familydefault}{\sfdefault} 

% --- Colori e Grafica ---
\usepackage{xcolor}
\usepackage{tikz}
\usepackage{amssymb}
\usepackage{amsmath} 
\usepackage{amsthm}
\usepackage{mdframed}
\usepackage{cancel}
\usepackage{booktabs}
\usepackage{fancyhdr}
\usepackage{enumitem}
\usepackage{multicol} % Per layout a 2 colonne (Cheat Sheet)
\usepackage[most]{tcolorbox} 

\usetikzlibrary{shapes, arrows.meta, positioning, chains, fit, calc, matrix, trees, backgrounds}

% --- Definizione Colori (Cotton + Cherry Red Palette) ---
\definecolor{cotton}{RGB}{237, 235, 221}    % #EDEBDD
\definecolor{cherryRed}{RGB}{129, 1, 0}     % #810100
\definecolor{maroon}{RGB}{99, 0, 0}         % #630000
\definecolor{noirBlack}{RGB}{27, 23, 23}    % #1B1717

% Mappatura sui colori logici del progetto
\colorlet{mainColor}{cherryRed}
\colorlet{accentColor}{maroon}
\colorlet{lightBg}{cotton}
\colorlet{textColor}{noirBlack}

% --- Header & Footer Setup ---
\pagestyle{fancy}
\fancyhf{} 
\renewcommand{\headrulewidth}{1.5pt}
\renewcommand{\headrule}{\hbox to\headwidth{\color{mainColor}\leaders\hrule height \headrulewidth\hfill}}

\lhead{\textbf{\textcolor{textColor}{Sistemi Lineari e Matrici}}}
\cfoot{\thepage}

% --- COMANDO COPERTINA (Left Aligned as requested) ---
\newcommand{\customcover}[2]{
    \begin{titlepage}
    \thispagestyle{empty}
    \begin{tikzpicture}[remember picture, overlay]
        % 1. Sfondo Laterale (Scuro - Noir Black)
        \fill[noirBlack] (current page.north west) rectangle ($(current page.south west) + (4cm,0)$);
        
        % 2. Linea Accento (Cherry Red)
        \draw[cherryRed, line width=4pt] ($(current page.north west) + (4.2cm,-2cm)$) -- ($(current page.south west) + (4.2cm,2cm)$);
        
        % 3. Contenuto Testuale (Allineato a SINISTRA, vicino alla fascia scura)
        % Anchor North West, posizionato rispetto al bordo sinistro (dopo la fascia nera)
        \node[anchor=north west, inner sep=0, text width=13cm, align=left] at ($(current page.north west) + (5cm,-5cm)$) {
            
                {\fontsize{18}{22}\selectfont \textcolor{maroon}{Università di Pisa}}\\[0.2cm]
                {\fontsize{16}{20}\selectfont \textcolor{maroon}{Corso di Laurea in Informatica}}\\[2cm]
                
                {\fontsize{60}{70}\selectfont \textbf{\textcolor{cherryRed}{LA BIBBIA}}}\\[0.5cm]
                
                
                {\fontsize{35}{45}\selectfont \textbf{\textcolor{noirBlack}{ di ALGEBRA LINEARE}}}\\[1.5cm]
                
                % Linea divisoria CherryRed
                \textcolor{cherryRed}{\rule{13cm}{4pt}}\\[1cm]
                
                {\Large \textbf{#1}}\\[0.2cm]
                {\large \textit{#2}}\\[1.5cm]
                
                {\large Docente titolare:}\\[0.2cm]
                {\Large \textbf{Prof.ssa Silvia Benvenuti}}\\[2cm]
                
                \textbf{\large Autore del riassunto:}\\
                {\large \textbf{Joseph Zucchelli (ilbroccolo)}}\\[2cm]
                
                \vfill
                {\large \textbf{Anno Accademico 2025 - 2026}}
            
        };
        
        % Badge (In basso a destra)
        \node[anchor=south east] at ($(current page.south east) + (-2cm, 2cm)$) {
            \tikz\draw[fill=cherryRed, draw=none] (0,0) rectangle (2,2) node[midway, white, font=\bfseries\huge] {AL};
        };

    \end{tikzpicture}
    \end{titlepage}
    \clearpage
}

% --- Setup Titoli ---
\usepackage{titlesec}
\titleformat{\section}
{\color{mainColor}\normalfont\Large\bfseries}
{\color{mainColor}\thesection}{1em}{}[\color{mainColor}\hrule height 1pt]

\titleformat{\subsection}
{\color{mainColor}\normalfont\large\bfseries}
{\color{mainColor}\thesubsection}{1em}{}

% --- Box Styles (Themed) ---

% Concept Box (Using lighter variant of mainColor)
\newtcolorbox{concept}[1]{
  colback=lightBg,
  colframe=accentColor,
  title={#1},
  coltitle=white,
  fonttitle=\bfseries,
  sharp corners=downhill,
  enhanced,
  drop shadow
}

% Definition Box (Replacing previous style)
\newtcolorbox{definitionbox}{
  colback=white,
  colframe=cherryRed,
  fonttitle=\bfseries,
  coltitle=white,
  title={Definizione},
  sharp corners,
  enhanced,
  borderline west={4pt}{0pt}{cherryRed}
}

% --- BOX PERSONALIZZATI ---

% Esempio (Redefinition using theme colors)
\newtcbtheorem[number within=section]{example}{Esempio}{
    enhanced,
    colback=white,
    colframe=accentColor,
    coltitle=white,
    fonttitle=\bfseries,
    attach boxed title to top left={yshift=-2mm, xshift=2mm},
    boxed title style={colback=cherryRed, sharp corners},
    breakable
}{ex}

% Nota Manoscritta
\newtcbtheorem[number within=section]{note}{Nota Manoscritta}{
    enhanced,
    frame hidden,
    colback=gray!10,
    borderline west={2mm}{0pt}{accentColor},
    fonttitle=\bfseries\color{maroon},
    coltitle=maroon,
    title={#1},
    breakable
}{note}

% Definizione (Lezione 6 style adapted to theme)
\newtcbtheorem[number within=section]{definition}{Definizione}{
    enhanced,
    colback=lightBg,
    colframe=cherryRed,
    coltitle=white,
    fonttitle=\bfseries,
    attach boxed title to top left={yshift=-2mm, xshift=2mm},
    boxed title style={colback=maroon, sharp corners},
    breakable
}{def}

% Osservazione
\newtcbtheorem[number within=section]{observation}{Osservazione}{
    enhanced,
    frame hidden,
    colback=white,
    borderline west={1mm}{0pt}{cherryRed},
    fonttitle=\bfseries\color{cherryRed},
    coltitle=cherryRed,
    title={#1},
    breakable
}{obs}

% Proposizione (New Lezione 5)
\newtcbtheorem[number within=section]{proposition}{Proposizione}{
    enhanced,
    colback=white,
    colframe=accentColor,
    coltitle=white,
    fonttitle=\bfseries,
    attach boxed title to top left={yshift=-2mm, xshift=2mm},
    boxed title style={colback=cherryRed, sharp corners},
    breakable
}{prop}

\begin{document}

% Copertina Aggiornata
\customcover{Appunti Completi: Sistemi e Matrici}{Trascrizione analitica Lezioni 1, 2, 3, 4, 5, 6}

\tableofcontents
\newpage

% Include Chapters

\section{Matrici}
\emph{[Rif. Lezioni 2 e 3]}

\subsection{Definizioni Fondamentali}
Per studiare i sistemi lineari, basta studiare le loro matrici associate.

\subsubsection{Matrice Associata e Completa}
Dato un sistema generico, definiamo:
\begin{itemize}
    \item \textbf{Matrice Associata ($A$):} Tabella $m \times n$ dei coefficienti. $A = (a_{ij})$.
    \item \textbf{Vettore Termini Noti ($b$):} Matrice colonna $m \times 1$.
    \item \textbf{Matrice Completa ($A|b$):} Matrice $m \times (n+1)$ ottenuta affiancando $b$ ad $A$.
\end{itemize}

\subsection{Matrici Quadrate e Tipologie Speciali}
Se $m=n$, la matrice è \textbf{Quadrata} di \textbf{ORDINE} $n$.

\begin{definitionbox}
\textbf{1. Diagonale Principale:} Insieme degli elementi $a_{ij}$ con $i=j$ ($a_{11}, a_{22}, \dots$).

\textbf{2. Matrice Diagonale:} Se $a_{ij} = 0$ per ogni $i \neq j$.
$$ \text{Esempio: } A = \begin{pmatrix} 1 & 0 & 0 \\ 0 & 3 & 0 \\ 0 & 0 & 5 \end{pmatrix} $$

\textbf{3. Matrice Triangolare Superiore:} Se tutti gli elementi \textbf{SOTTO} la diagonale principale sono 0 ($a_{ij}=0 \ \forall i > j$).
$$ \text{Esempio: } A = \begin{pmatrix} 1 & 3 & 5 \\ 0 & 2 & 4 \\ 0 & 0 & 6 \end{pmatrix} $$

\textbf{4. Matrice Triangolare Inferiore:} Se tutti gli elementi \textbf{SOPRA} la diagonale principale sono 0 ($a_{ij}=0 \ \forall i < j$).
$$ \text{Esempio: } A = \begin{pmatrix} 1 & 0 \\ 2 & 3 \end{pmatrix} $$

\textbf{Osservazione:} Una matrice diagonale è contemporaneamente triangolare superiore e inferiore.
\end{definitionbox}

\newpage
\section{Sistemi Lineari e Gauss}
\emph{[Rif. Lezione 1, Lezione 2 e 3]}

\subsection{Concetti Introduttivi (Lezione 1)}

\subsubsection{Premesse Fondamentali}
Salvo esplicito avviso contrario, supponiamo sempre di lavorare in $\mathbb{R}$ (insieme dei numeri reali).

\subsubsection{Equazioni Lineari}
\textbf{Definizione:} Un'equazione lineare è un'equazione di I grado (in cui l'incognita compare al grado 1).

\paragraph{Caso 1: Equazioni in una variabile}
Forma normale: $ax + b = 0$.
\begin{itemize}
    \item Se $a \neq 0 \implies$ \textbf{1 soluzione}: $x = -\frac{b}{a}$.
    \item Se $a = 0$ e $b \neq 0 \implies$ \textbf{Impossibile}.
    \item Se $a = 0$ e $b = 0 \implies$ \textbf{Indeterminata} ($\forall x \in \mathbb{R}$).
\end{itemize}

\paragraph{Caso 2: Equazioni in due variabili}
Esempio: $x + y = 1$.
Soluzioni: $\{ (t, 1-t) \mid t \in \mathbb{R} \}$.
\textbf{Interpretazione Geometrica:} I punti di una retta.

\subsection{Sistemi di Equazioni Lineari}
\textbf{Definizione:} Insieme di equazioni che devono essere soddisfatte \textbf{contemporaneamente}.

\subsubsection{Interpretazione Geometrica (Sistemi $2 \times 2$)}
Cercare l'intersezione tra due rette.
\begin{center}
\begin{tabular}{ccc}
\toprule
\textbf{Configurazione} & \textbf{Soluzioni} & \textbf{Esempio} \\
\midrule
\textbf{Rette Incidenti} & \textbf{1 Soluzione} & $\begin{cases} x+y=1 \\ 2x-y=0 \end{cases}$ \\
\midrule
\textbf{Rette Parallele} & \textbf{0 Soluzioni} & $\begin{cases} x+y=1 \\ x+y=2 \end{cases}$ \\
\midrule
\textbf{Rette Coincidenti} & \textbf{$\infty$ Soluzioni} & $\begin{cases} x+y=1 \\ 2x+2y=2 \end{cases}$ \\
\bottomrule
\end{tabular}
\end{center}

\subsection{Generalizzazione}
Un sistema lineare può essere:
\begin{enumerate}
    \item \textbf{Possibile (Compatibile):}
    \begin{itemize}
        \item 1 Soluzione unica.
        \item $\infty$ Soluzioni.
    \end{itemize}
    \item \textbf{Impossibile:} 0 Soluzioni.
\end{enumerate}

\section{Algoritmo di Eliminazione di Gauss}

\subsection{Sistemi Equivalenti}
Due sistemi si dicono \textbf{EQUIVALENTI} se hanno lo stesso insieme di soluzioni.

\begin{concept}{Sistemi Triangolari}
I sistemi triangolari superiori sono facili da risolvere con la \textbf{Sostituzione a Ritroso}.
\end{concept}

\subsection{Le 3 Mosse di Gauss}
Operazioni lecite sulle righe della matrice completa $(A|b)$:
\begin{enumerate}
    \item[(I)] Scambiare due righe ($R_i \leftrightarrow R_j$).
    \item[(II)] Moltiplicare una riga per uno scalare $\lambda \neq 0$ ($R_i \to \lambda R_i$).
    \item[(III)] Sommare a una riga un multiplo di un'altra ($R_i \to R_i + k R_j$).
\end{enumerate}

\subsection{Esempi Pratici}

\subsubsection{Esempio 1: Sistema $3 \times 3$ Completo}
$$ \begin{cases} x+y+z=6 \\ 2x+y-z=3 \\ -x-y+2z=5 \end{cases} \implies C = \left(\begin{array}{ccc|c} 1 & 1 & 1 & 6 \\ 2 & 1 & -1 & 3 \\ -1 & -1 & 2 & 5 \end{array}\right) $$
\textbf{Passo 1:} Uccidere i termini sotto il primo pivot.
\begin{itemize}
    \item $R_2 \to R_2 - 2R_1$: $(0, -1, -3, -9)$
    \item $R_3 \to R_3 + R_1$: $(0, 0, 3, 11)$
\end{itemize}
$$ \left(\begin{array}{ccc|c}  1 & 1 & 1 & 6 \\  0 & -1 & -3 & -9 \\  0 & 0 & 3 & 11  \end{array}\right) $$
Risolvendo a ritroso: $z = 11/3$, $y = -2$, $x = 13/3$.

\subsubsection{Esempio 2: Impossibile}
$$ \left(\begin{array}{cc|c} 1 & 1 & 3 \\ 0 & -4 & -13 \end{array}\right) $$
Se avessimo ottenuto una riga $\begin{pmatrix} 0 & 0 & | & k \end{pmatrix}$ con $k \neq 0$, il sistema sarebbe impossibile.
Esempio:
$$ \left(\begin{array}{ccc|c} 2 & \dots & \dots & \dots \\ 0 & 0 & 0 & 2 \\ 0 & 0 & 0 & 0 \end{array}\right) \implies \text{IMPOSSIBILE} $$

\subsubsection{Esempio 3: Famiglia di Soluzioni}
$$ \left(\begin{array}{ccc|c} 1 & 1 & 1 & 3 \\ 0 & -1 & 1 & 1 \end{array}\right) $$
Ponendo $z=t$ (libero): $y = t-1$, $x = 4-2t$.
Soluzioni: $(4-2t, t-1, t)$ per ogni $t \in \mathbb{R}$.

\subsection{Riepilogo e Morale}
\begin{itemize}
    \item Riga $0 \dots 0 | b \ (b \neq 0) \implies$ \textbf{0 soluzioni}.
    \item Riga $0=0 \implies$ Ridondante.
    \item Sistema quadrato triangolare con pivot $\neq 0 \implies$ \textbf{1 soluzione}.
\end{itemize}

\newpage
\newgeometry{top=2cm, bottom=2cm, left=1.8cm, right=1.8cm}
\setlength{\parindent}{0pt}
\setlength{\parskip}{6pt}

% --- INTESTAZIONE ---
\begin{center}
    {\Huge \textbf{\color{cherryRed} Lezione 4: Algoritmo di Gauss}} \\
    \vspace{0.3cm}
    {\Large \textit{Trascrizione Integrale del PDF "lezione 4\_090226"}}
    \vspace{0.5cm}
    \hrule height 1pt
\end{center}

% =========================================================================
% PAGINA 1
% =========================================================================
\section{Introduzione e Risoluzione a Ritroso (Pag. 1)}

\begin{example}{Esempio 2}{}
Adesso la matrice è triangolare superiore ($\Rightarrow$ "Tri Sup"). Risolvo il sistema "a ritroso".

Sistema:
\[
\begin{cases}
x + y + z = 6 \\
y - 3z = -9 \\
7z = 17
\end{cases}
\]

\textbf{Svolgimento:}
1. Dalla terza equazione: $z = \frac{17}{7}$.
2. Sostituisco nella seconda: $y = -9 + 3z = -9 + 3\left(\frac{17}{7}\right) = -9 + \frac{51}{7} = \frac{-63+51}{7} = -\frac{12}{7}$.
3. Sostituisco nella prima: $x = 6 - y - z = 6 - \left(-\frac{12}{7}\right) - \frac{17}{7} = 6 + \frac{12}{7} - \frac{17}{7} = 6 - \frac{5}{7} = \frac{42-5}{7} = \frac{37}{7}$.

\textbf{Nota:} Per avere la "Tri Sup" basta scambiare le righe $C_2 \leftrightarrow C_3$ se necessario.
\end{example}

\begin{example}{Esempio 3 (Veloce)}{}
\[
\begin{pmatrix} 1 & 1 \\ 0 & 1 \end{pmatrix} \implies \begin{cases} y=1 \\ x+y=2 \implies x=1 \end{cases}
\]
Soluzione: $(1, 1)$.
\end{example}

\begin{example}{Esempio 4 (Impossibile)}{}
\[
\begin{cases}
x+y=2 \\
x+y=3 \\
2x+2y=4
\end{cases}
\]
\begin{note}{Osservazione}{}
Si vede "a occhio" che la $1^a$ e la $2^a$ equazione sono incompatibili $\implies$ il sistema è impossibile.
Ma supponiamo che io non me ne accorga.
\end{note}

Matrice associata:
\[
\begin{pmatrix}
1 & 1 & | & 2 \\
1 & 1 & | & 3 \\
2 & 2 & | & 4
\end{pmatrix}
\]
Ammazzo i pivot sotto $A_{11}$:
\begin{itemize}
    \item $C_2 \to C_2 - C_1$
    \item $C_3 \to C_3 - 2C_1$
\end{itemize}
Otteniamo:
\[
\begin{pmatrix}
1 & 1 & | & 2 \\
0 & 0 & | & 1 \\
0 & 0 & | & 0
\end{pmatrix}
\]
La seconda riga corrisponde all'equazione $0x + 0y = 1 \implies 0=1$. \textbf{IMPOSSIBILE}.
\end{example}

% =========================================================================
% PAGINA 2
% =========================================================================
\section{Variabili Libere e Interpretazione Geometrica (Pag. 2)}

\begin{example}{Esempio 5}{}
\[
\begin{cases}
x+y+z=3 \\
2x+y+3z=7
\end{cases}
\xrightarrow{\text{Matrice}}
\begin{pmatrix} 1 & 1 & 3 & | & 3 \\ 2 & 1 & 3 & | & 7 \end{pmatrix}
\]
Applicazione Gauss: $C_2 \to C_2 - 2C_1$.
\[
\begin{pmatrix} 1 & 1 & 3 & | & 3 \\ 0 & -1 & -1 & | & 1 \end{pmatrix}
\]
\end{example}

\begin{note}{Morale}{}
La forma a scalini fa emergere subito una "famiglia" di soluzioni, in questo caso dipendenti da 1 parametro (gradi di libertà).
Chiamiamo \textbf{Variabili Libere} quelle corrispondenti alle colonne in cui NON c'è il pivot.
\end{note}

Nel sistema ridotto:
\[
\begin{cases}
x+y+z=3 \\
-y-z=1
\end{cases}
\]
Variabile libera: $z$ (terza colonna senza pivot). Poniamo $z = t$.
Svolgimento a ritroso:
1. $-y - t = 1 \implies y = -1 - t$.
2. $x + (-1-t) + t = 3 \implies x - 1 = 3 \implies x = 4$.

\textbf{Soluzione:} $(4, -1-t, t)$.
Famiglia di soluzioni $\infty^1$.

\textbf{Riassumendo quanto visto dagli esempi:}
\begin{itemize}
    \item Se compare una riga $0 \ 0 \ \dots \ 0 \ | \ b$ con $b \neq 0 \implies$ \textbf{Nessuna soluzione}.
    \item Una o più righe tutte di zeri $\implies$ qualche equazione ridondante (si possono cancellare).
    \item Sistema quadrato che diventa triangolare con pivot tutti non nulli $\implies$ \textbf{1! Soluzione}.
\end{itemize}

\begin{note}{Importante}{}
Il numero dei pivot misura quanta "info indipendente" c'è.
\end{note}

% =========================================================================
% PAGINA 3
% =========================================================================
\section{Esempi Veloci e Introduzione Teoria (Pag. 3)}

\begin{itemize}
    \item $\begin{cases} x+y=4 \\ 2x+4y=8 \end{cases} \implies \infty^1$ sol.
    \item $\begin{cases} 2x+2y=4 \\ 2x+2y=9 \end{cases} \implies 0$ sol (rette parallele distinte).
    \item $\begin{cases} 2x+y=5 \\ 2x-y=1 \end{cases} \implies 1!$ sol (rette incidenti).
\end{itemize}

A questo punto enunciamo la \textbf{Teoria Generale}: ALGORITMO DI GAUSS (ripetuto per comodità alla pagina successiva).

% =========================================================================
% PAGINA 4
% =========================================================================
\section{Algoritmo di Gauss e Mosse Elementari (Pag. 4)}

\subsection{Equivalenza di Sistemi}
Sia $C = (A|b)$ la matrice completa di un sistema lineare.
Le \textbf{Mosse di Gauss} (Operazioni Elementari) sono:
\begin{enumerate}
    \item[(I)] Scambiare tra loro due righe ($R_i \leftrightarrow R_j$).
    \item[(II)] Moltiplicare una riga per una costante $\lambda \neq 0$.
    \item[(III)] Sostituire una riga $C_i$ con la somma di $C_i$ con un multiplo di un'altra riga: $C_i \to C_i + \lambda C_j$.
\end{enumerate}

Due sistemi ottenuti l'uno dall'altro tramite mosse di Gauss sono tra loro \textbf{EQUIVALENTI}, cioè hanno le stesse soluzioni.

\subsection{Metodo di Eliminazione di Gauss}
È una procedura (algoritmo) che consente di trasformare un qualsiasi sistema lineare in un sistema a scalini equivalente.
\textbf{Obiettivo:} Ottenere $C_{ij} = 0$ sotto la diagonale.

% =========================================================================
% PAGINA 5
% =========================================================================
\section{Procedimento Passo-Passo (Pag. 5)}

\textbf{Passo 1:}
Considera la prima colonna di $C$, $C_1$.
\begin{itemize}
    \item Se contiene solo zeri, lasciala stare e ricomincia dalla colonna 2.
    \item Se contiene qualche elemento $\neq 0$, a meno di scambiare tra loro le righe, portati nella situazione $a_{11} \neq 0$ (questo sarà il \textbf{Primo Pivot} $P_1$).
    \item Ora "AMMAZZA" tutti i coefficienti sotto $P_1$. Per farlo, basta sostituire $C_i \to C_i - \frac{C_{i1}}{P_1} C_1$.
\end{itemize}
Ti sei ricondotto a una matrice di forma:
\[
\begin{pmatrix}
P_1 & \times & \dots \\
0 & \times & \dots \\
0 & \times & \dots
\end{pmatrix}
\]

\textbf{Passo 2:}
Ripeti (I) e (II) sulla sottomatrice $(A'|b')$.
E così via fino alla fine.

% =========================================================================
% PAGINA 6
% =========================================================================
\section{Variabili Libere e Pivot (Pag. 6)}

Termini con una situazione di questo tipo (Matrice a scalini):
\[
\begin{pmatrix}
P_1 & \times & \times & \dots & | & \times \\
0 & P_2 & \times & \dots & | & \times \\
0 & 0 & 0 & P_3 & \dots & | & \times \\
\vdots & & & & & & \vdots
\end{pmatrix}
\]
Notiamo colonne senza pivot (es. colonna 3).
\begin{note}{Definizione}{}
Chiamo \textbf{VARIABILI LIBERE} quelle corrispondenti alle colonne in cui NON ho pivot.
Sono $n-r$ (dove $r$ è il numero di pivot).
\end{note}

Il processo di eliminazione ha termine quando $r=m$ oppure, se $r<m$, le ultime $m-r$ righe sono tutte nulle (nella matrice dei coefficienti).
Se $r=m=n$, la matrice è triangolare superiore con pivot sulla diagonale principale.

% =========================================================================
% PAGINA 7
% =========================================================================
\section{Esempio di Applicazione (Pag. 7)}

Matrice $C$:
\[
C = \begin{pmatrix}
0 & 1 & 1 & | & 0 \\
1 & 1 & 2 & | & -3 \\
-1 & 2 & 1 & | & 1
\end{pmatrix}
\]
1. Scambio $R_1 \leftrightarrow R_2$ (per avere pivot in alto a sinistra):
\[
\begin{pmatrix}
1 & 1 & 2 & | & -3 \\
0 & 1 & 1 & | & 0 \\
-1 & 2 & 1 & | & 1
\end{pmatrix}
\]
2. $C_3 \to C_3 + C_1$:
\[
\begin{pmatrix}
1 & 1 & 2 & | & -3 \\
0 & 1 & 1 & | & 0 \\
0 & 3 & 3 & | & -2
\end{pmatrix}
\]
3. $C_3 \to C_3 - 3C_2$:
\[
\begin{pmatrix}
1 & 1 & 2 & | & -3 \\
0 & 1 & 1 & | & 0 \\
0 & 0 & 0 & | & -2
\end{pmatrix}
\]
L'ultima riga è $0 = -2 \implies$ \textbf{IMPOSSIBILE}.

\subsection{Algoritmo di Gauss-Jordan}
Oltre ad arrivare a una matrice a scalini, posso renderla di una forma particolare (Semplice):
1. Tutti i pivot uguali a 1 (basta dividere la riga per il pivot).
2. Tutti i coefficienti \textbf{SOPRA} i pivot uguali a 0.
Come si fa? Si opera "dal basso" verso l'alto.

Questa si chiama \textbf{MATRICE COMPLETAMENTE RIDOTTA DEL SISTEMA}.

% =========================================================================
% PAGINA 8
% =========================================================================
\section{Esercizio Svolto (Pag. 8)}

\textbf{Testo:}
Dato il sistema:
\[
\begin{cases}
x + 3y + z - w = 2 \\
3x + 9y + 4z + w = 1 \\
2x + y + 5z + 2w = 0 \\
y - z - w = 2
\end{cases}
\]
1. Scrivere la matrice completa.
2. Applicando Gauss, trasformare in matrice a scalini (Matrice Ridotta).
3. Determinare la matrice COMPLETAMENTE RIDOTTA.
4. Cosa puoi osservare sulle soluzioni?

\textbf{Svolgimento:}
Matrice $C$:
\[
\begin{pmatrix}
1 & 3 & 1 & -1 & | & 2 \\
3 & 9 & 4 & 1 & | & 1 \\
2 & 1 & 5 & 2 & | & 0 \\
0 & 1 & -1 & -1 & | & 2
\end{pmatrix}
\]
Passaggio 1: Eliminare sotto il primo pivot (1).
$C_2 \to C_2 - 3C_1$ e $C_3 \to C_3 - 2C_1$.
\[
\begin{pmatrix}
1 & 3 & 1 & -1 & | & 2 \\
0 & 0 & 1 & 4 & | & -5 \\
0 & -5 & 3 & 4 & | & -4 \\
0 & 1 & -1 & -1 & | & 2
\end{pmatrix}
\]
Nota: Sotto il secondo elemento della prima riga c'è uno zero, dobbiamo scambiare righe per avere un pivot. Scambio $C_2 \leftrightarrow C_4$ (o $C_2 \leftrightarrow C_3$). Scambiamo con l'ultima per comodità (come nell'appunto):
\[
\begin{pmatrix}
1 & 3 & 1 & -1 & | & 2 \\
0 & 1 & -1 & -1 & | & 2 \\
0 & -5 & 3 & 4 & | & -4 \\
0 & 0 & 1 & 4 & | & -5
\end{pmatrix}
\]
Continuo eliminazione: $C_3 \to C_3 + 5C_2$.
\[
\begin{pmatrix}
1 & 3 & 1 & -1 & | & 2 \\
0 & 1 & -1 & -1 & | & 2 \\
0 & 0 & -2 & -1 & | & 6 \\
0 & 0 & 1 & 4 & | & -5
\end{pmatrix}
\]
Scambio $C_3 \leftrightarrow C_4$ per avere 1 come pivot (più comodo) oppure procedo. Nell'appunto si vede che si arriva a:
\[
\begin{pmatrix}
1 & 3 & 1 & -1 & | & 2 \\
0 & 1 & -1 & -1 & | & 2 \\
0 & 0 & 1 & 4 & | & -5 \\
0 & 0 & -2 & -1 & | & 6
\end{pmatrix}
\]
Ultimo passaggio: $C_4 \to C_4 + 2C_3$.
\[
\begin{pmatrix}
1 & 3 & 1 & -1 & | & 2 \\
0 & 1 & -1 & -1 & | & 2 \\
0 & 0 & 1 & 4 & | & -5 \\
0 & 0 & 0 & 7 & | & -4
\end{pmatrix}
\]
Matrice a scalini (anzi Triangolare Superiore).
Ci fermiamo.
$r = \# \text{pivot} = 4$.

% =========================================================================
% PAGINA 9
% =========================================================================
\section{Continuazione Esercizio (Pag. 9)}

La matrice ridotta mi fornisce automaticamente la soluzione (che è unica, $r=n=4$).
Dall'ultima riga:
\[ 7w = -4 \implies w = -4/7 \]
Ora usiamo Gauss-Jordan (uccidere i coefficienti sopra).
Normalizzo l'ultimo pivot ($P_4=7 \to 1$): $C_4 \to \frac{1}{7}C_4$.
\[
\begin{pmatrix}
1 & 3 & 1 & -1 & | & 2 \\
0 & 1 & -1 & -1 & | & 2 \\
0 & 0 & 1 & 4 & | & -5 \\
0 & 0 & 0 & 1 & | & -4/7
\end{pmatrix}
\]
Operazioni a ritroso:
\begin{itemize}
    \item $C_3 \to C_3 - 4C_4$
    \item $C_2 \to C_2 + C_4$
    \item $C_1 \to C_1 + C_4$
\end{itemize}
Matrice diventa:
\[
\begin{pmatrix}
1 & 3 & 1 & 0 & | & 2 - 4/7 = 10/7 \\
0 & 1 & -1 & 0 & | & 2 - 4/7 = 10/7 \\
0 & 0 & 1 & 0 & | & -5 + 16/7 = -19/7 \\
0 & 0 & 0 & 1 & | & -4/7
\end{pmatrix}
\]
Proseguo eliminando sopra il pivot della riga 3 ($P_3=1$):
\begin{itemize}
    \item $C_2 \to C_2 + C_3$
    \item $C_1 \to C_1 - C_3$
\end{itemize}
\[
\begin{pmatrix}
1 & 3 & 0 & 0 & | & 10/7 - (-19/7) = 29/7 \\
0 & 1 & 0 & 0 & | & 10/7 - 19/7 = -9/7 \\
0 & 0 & 1 & 0 & | & -19/7 \\
0 & 0 & 0 & 1 & | & -4/7
\end{pmatrix}
\]
Infine elimino sopra $P_2$: $C_1 \to C_1 - 3C_2$.
\[
\begin{pmatrix}
1 & 0 & 0 & 0 & | & 29/7 - 3(-9/7) = 56/7 = 8 \\
0 & 1 & 0 & 0 & | & -9/7 \\
0 & 0 & 1 & 0 & | & -30/7 \text{ (Nota: correzione calcolo a margine)}\\
0 & 0 & 0 & 1 & | & -4/7
\end{pmatrix}
\]
(Nota: ci sono correzioni manuali nei calcoli sul foglio, ho riportato i valori finali visibili: $x=8$ o valori simili corretti a penna come $30/7$).

% =========================================================================
% PAGINA 10
% =========================================================================
\section{Teorema Sistema Quadrato (Pag. 10)}

\begin{tcolorbox}[colback=white, colframe=black, title=TEOREMA]
Un sistema lineare quadrato ammette 1! soluzione $\iff$ i pivot della sua matrice ridotta sono tutti $\neq 0$ (cioè sono $n$).
In tal caso la matrice ridotta è triangolare superiore.
\end{tcolorbox}

\begin{note}{Osservazione}{}
Come abbiamo ben visto dagli esempi, i pivot non sono univocamente determinati, perché dipendono dalle scelte che effettuiamo nell'applicare l'algoritmo (scambi di riga).
Tuttavia, si può dimostrare che: il NUMERO dei pivot è COSTANTE.
\end{note}

\textbf{Proposizione:} Se una matrice ammette $r$ pivots per una serie di mosse, ne ammette $r$ per un'altra qualsiasi scelta.
TRADUZIONE: I pivot non sono univoci, ma il loro NUMERO sì!

% =========================================================================
% PAGINA 11
% =========================================================================
\section{Rango e Rouché-Capelli (Pag. 11)}

\begin{note}{Definizione}{}
Il numero dei pivot di una qualsiasi riduzione a scale di una matrice $A$ si dice \textbf{RANGO} della matrice $A$.
Si scrive $r = \text{rk}(A)$.
Vale sempre $r \le \min(m, n)$.
\end{note}
Una matrice quadrata di ordine $n$ si dice \textbf{NON SINGOLARE} se $\text{rk}(A)=n$, SINGOLARE altrimenti.

\begin{tcolorbox}[colback=lightBg, colframe=cherryRed, title=TEOREMA DI ROUCHÉ-CAPELLI]
Sia $C=(A|b)$ la matrice completa di un sistema lineare $m \times n$.
Supponiamo di aver ridotto a scale la matrice $A$.
\begin{itemize}
    \item Il sistema è \textbf{COMPATIBILE} (ha soluzioni) $\iff \text{rk}(A) = \text{rk}(A|b)$.
    Ovvero non compare una riga del tipo $(0 \dots 0 | d)$ con $d \neq 0$.
    \item Se compatibile:
    \begin{itemize}
        \item Ammette \textbf{1! soluzione} se $r = n$.
        \item Ammette \textbf{infinite soluzioni} dipendenti da $n-r$ parametri se $r < n$.
    \end{itemize}
\end{itemize}
\end{tcolorbox}

% =========================================================================
% PAGINA 12
% =========================================================================
\section{Dimostrazione Costruttiva (Pag. 12)}

Dimostrazione "costruttiva" usando la forma ridotta.
Se la matrice $(A|b)$ ridotta è della forma:
\[
\begin{pmatrix}
\dots & | & d_1 \\
\vdots & | & \vdots \\
0 \dots 0 & | & d_r \\
0 \dots 0 & | & d_{r+1} \\
\vdots & | & \vdots \\
0 \dots 0 & | & d_m
\end{pmatrix}
\]
Le righe $r+1, \dots, m$ corrispondono alle equazioni $0 = d_{r+1}, \dots, 0 = d_m$.
\begin{itemize}
    \item Se uno dei $d_i$ con $i > r$ fosse $\neq 0 \implies$ Impossibile.
    \item Se tutti i $d_i$ sono $0 \implies$ Posso dimenticarmi delle ultime equazioni (ridondanti).
\end{itemize}
Rimane un sistema $r \times n$.
Chiamo \textbf{variabili libere} quelle corrispondenti alle colonne senza pivot. Sono $n-r$. A ciascuna assegno un parametro $t_j$.
Mi riconduco a un sistema triangolare risolvibile.

% =========================================================================
% PAGINA 13
% =========================================================================
\section{Esempio Dimostrativo (Pag. 13)}

Consideriamo il sistema lineare:
\[
\begin{cases}
2x_2 + 2x_3 - 3x_4 - x_5 = 0 \\
x_1 - 4x_2 + 2x_3 - 3x_4 + x_5 = 1 \\
x_4 + x_5 = 3
\end{cases}
\]
Matrice:
\[
C = \begin{pmatrix}
0 & 2 & 2 & -3 & -1 & | & 0 \\
1 & -4 & 2 & -3 & 1 & | & 1 \\
0 & 0 & 0 & 1 & 1 & | & 3
\end{pmatrix}
\]
Scambio $R_1 \leftrightarrow R_2$:
\[
\begin{pmatrix}
1 & -4 & 2 & -3 & 1 & | & 1 \\
0 & 2 & 2 & -3 & -1 & | & 0 \\
0 & 0 & 0 & 1 & 1 & | & 3
\end{pmatrix}
\]
È già a scalini!
$r = \#\text{pivot} = 3$ (Pivot: 1, 2, 1).
Sistema compatibile (non ci sono righe impossibili).
Numero incognite $n=5$.
Infinite soluzioni dipendenti da $n-r = 5-3 = 2$ parametri.

Le colonne senza pivot sono la 3 e la 5 ($x_3, x_5$).
Pongo $x_3 = t_3, x_5 = t_5$.
Risolvo:
\begin{itemize}
    \item $x_4 + x_5 = 3 \implies x_4 = 3 - t_5$
    \item $2x_2 + 2t_3 - 3(3-t_5) - t_5 = 0 \implies$ trovo $x_2$.
    \item $x_1 - \dots = 1 \implies$ trovo $x_1$.
\end{itemize}

% =========================================================================
% PAGINA 14 (ESERCIZIO PARAMETRICO)
% =========================================================================
\section{Esercizio Parametrico (Pag. 14)}

\textbf{Esercizio:} Determinare i valori di $a$ per i quali il seguente sistema presenta:
i) Nessuna soluzione; ii) 1! soluzione; iii) Più di una soluzione.

\[
\begin{cases}
x + y - z = 1 \\
2x + 3y + az = 3 \\
x + 2y + 3z = 2
\end{cases}
\xrightarrow{\text{Matrice}}
\begin{pmatrix}
1 & 1 & -1 & | & 1 \\
2 & 3 & a & | & 3 \\
1 & a & 3 & | & 2
\end{pmatrix}
\]
Effettuo la riduzione a scalini:
1. $R_2 \to R_2 - 2R_1$.
2. $R_3 \to R_3 - R_1$.
\[
\begin{pmatrix}
1 & 1 & -1 & | & 1 \\
0 & 1 & a+2 & | & 1 \\
0 & a-1 & 4 & | & 1
\end{pmatrix}
\]
3. $R_3 \to R_3 - (a-1)R_2$.
Calcolo elemento $(3,3)$: $4 - (a-1)(a+2) = 4 - (a^2+2a-a-2) = 4 - (a^2+a-2) = 6 - a - a^2$.
Calcolo termine noto $(3,4)$: $1 - (a-1)\cdot 1 = 2 - a$.

Matrice ridotta:
\[
\begin{pmatrix}
1 & 1 & -1 & | & 1 \\
0 & 1 & a+2 & | & 1 \\
0 & 0 & 6-a-a^2 & | & 2-a
\end{pmatrix}
\]
Nota che $6-a-a^2 = -(a^2+a-6) = -(a+3)(a-2)$.

\textbf{Discussione:}
\begin{enumerate}
    \item[(i)] \textbf{NESSUNA SOLUZIONE:}
    Se il pivot è $0$ ma il termine noto è $\neq 0$.
    $6-a-a^2 = 0 \iff a=2 \lor a=-3$.
    \begin{itemize}
        \item Se $a=-3$: L'ultima riga diventa $0 \ | \ 2-(-3) = 5$. $0=5$ IMPOSSIBILE.
    \end{itemize}
    
    \item[(ii)] \textbf{1! SOLUZIONE:}
    Se il pivot è $\neq 0$.
    $6-a-a^2 \neq 0 \implies a \neq 2 \land a \neq -3$.
    In questo caso $r=n=3$.
    
    \item[(iii)] \textbf{PIÙ DI UNA SOLUZIONE (Infinite):}
    Se l'ultima riga è $0=0$.
    Serve $a=2$.
    Termine noto: $2-2=0$.
    Pivot: $0$.
    Riga: $0 \ 0 \ 0 \ | \ 0$.
    Il rango diventa 2 ($<3$). Infinite soluzioni ($\infty^1$).
\end{enumerate}

\restoregeometry

\newpage
\newgeometry{top=2.5cm, bottom=2.5cm, left=2cm, right=2cm}

% --- TITOLO ---
\begin{center}
    {\Huge \textbf{\color{cherryRed} Algebra Lineare e Geometria}} \\
    \vspace{0.5cm}
    {\Large \textit{Spazi Vettoriali e Sistemi Lineari (Lezione 5)}}
    \vspace{1cm}
    \hrule height 1pt
\end{center}

% -------------------------------------------------------------------------
\section{Spazi Vettoriali su $\mathbb{R}$}

\begin{concept}{Definizione di Spazio Vettoriale}{}
Uno \textbf{Spazio Vettoriale su $\mathbb{R}$} è un insieme $V$ dotato di due operazioni:
\begin{enumerate}
    \item \textbf{Somma:} $+ : V \times V \longrightarrow V$
    \item \textbf{Prodotto per scalare:} $\cdot : \mathbb{R} \times V \longrightarrow V$
\end{enumerate}
Tali operazioni devono soddisfare le seguenti 8 proprietà (assiomi).
\end{concept}

\subsection{Gli 8 Assiomi dello Spazio Vettoriale}

Le proprietà sono divise in due gruppi. Le prime 4 riguardano la struttura di gruppo additivo $(V, +)$.

\begin{proposition}{Proprietà della Somma (Gruppo Commutativo)}{}
Per ogni $u, v, w \in V$:
\begin{enumerate}
    \item \textbf{Associativa:} $(u + v) + w = u + (v + w)$
    \item \textbf{Commutativa:} $v + w = w + v$
    \item \textbf{Esistenza dell'Elemento Neutro:}
    \[ \exists 0 \in V : v + 0 = v, \quad \forall v \in V \]
    \item \textbf{Esistenza dell'Opposto:}
    \[ \forall v \in V, \exists -v \in V : v + (-v) = 0 \]
\end{enumerate}
\textbf{Nota:} L'insieme $(V, +)$ è dunque un \textit{Gruppo Commutativo}.
\end{proposition}

\begin{proposition}{Proprietà del Prodotto per Scalare}{}
Per ogni $v, w \in V$ e per ogni $\lambda, \mu \in \mathbb{R}$:
\begin{enumerate}[start=5]
    \item \textbf{Associatività (mista):} $(\lambda\mu) \cdot v = \lambda \cdot (\mu v)$
    \item \textbf{Elemento Neutro (moltiplicativo):} $1 \cdot v = v$
    \item \textbf{Distributiva (rispetto alla somma di vettori):} $\lambda(v + w) = \lambda v + \lambda w$
    \item \textbf{Distributiva (rispetto alla somma di scalari):} $(\lambda + \mu)v = \lambda v + \mu v$
\end{enumerate}
\end{proposition}

% -------------------------------------------------------------------------
\section{Lo Spazio $\mathbb{R}^n$}

\begin{concept}{Definizione di $\mathbb{R}^n$}{}
Definiamo l'insieme dei vettori colonna a $n$ componenti reali:
\[
\mathbb{R}^n = \left\{ v = \begin{pmatrix} x_1 \\ \vdots \\ x_n \end{pmatrix} : x_i \in \mathbb{R}, \forall i=1, \dots, n \right\}
\]
Esempio in $\mathbb{R}^3$:
\[
v = \begin{pmatrix} 1 \\ \pi \\ 7/2 \end{pmatrix} \in \mathbb{R}^3
\]
\end{concept}

\begin{proposition}{Operazioni in $\mathbb{R}^n$}{}
L'operazione $+$ in $\mathbb{R}^n$ gode delle seguenti proprietà (analogo generale dello spazio vettoriale):
\begin{itemize}
    \item $\forall v, w, z \in \mathbb{R}^n \implies v + (w + z) = (v + w) + z$ (Associativa)
    \item $\forall v, w \in \mathbb{R}^n \implies v + w = w + v$ (Commutativa)
    \item Elemento neutro: $0 = \begin{pmatrix} 0 \\ \vdots \\ 0 \end{pmatrix}$ tale che $v + 0 = v, \forall v$
    \item Opposto: $\forall v, \exists -v : v + (-v) = 0$
\end{itemize}
Inoltre valgono le proprietà legate agli scalari $\lambda, \mu \in \mathbb{R}$ (es. $(\lambda\mu)v = \lambda(\mu v)$).
\end{proposition}

% -------------------------------------------------------------------------
\section{Sistemi Lineari}

Un sistema lineare di $m$ equazioni in $n$ incognite ($m \times n$) si scrive come:
\[
\begin{cases}
a_{11}x_1 + \dots + a_{1n}x_n = b_1 \\
\vdots \\
a_{m1}x_1 + \dots + a_{mn}x_n = b_m
\end{cases}
\]
Dove $(x_1, \dots, x_n)^t$ sono le incognite in $\mathbb{R}$ e $b_i$ sono i termini noti.

\subsection{Matrici e Algoritmo di Gauss}

Consideriamo la matrice completa (o aumentata) del sistema $C = (A|b)$.
Applicando l'algoritmo di Gauss ($C \xrightarrow{\text{Gauss}}$), otteniamo una matrice a gradini (ridotta):

\[
\begin{pmatrix}
0 & 0 & p_1 & \times & \times & \dots & | & \times \\
0 & 0 & 0 & p_2 & \times & \dots & | & \times \\
\vdots & & & & & & & \vdots \\
0 & 0 & 0 & 0 & \dots & p_r & \dots & | & \times \\
0 & 0 & 0 & 0 & \dots & 0 & \dots & | & b'_{r+1} \\
\vdots & & & & & & & \vdots \\
0 & 0 & 0 & 0 & \dots & 0 & \dots & | & b'_m
\end{pmatrix}
\]

Dove:
\begin{itemize}
    \item $r$ è il numero di pivot (rango).
    \item Le prime $r$ righe contengono i pivot.
    \item Le ultime $m-r$ righe sono composte da zeri nella parte dei coefficienti.
\end{itemize}

\begin{concept}{Analisi delle Soluzioni (Rouché-Capelli)}{}
Sia $r$ il numero di pivot e $m$ il numero di equazioni.
\begin{enumerate}
    \item \textbf{Caso Impossibile:} Se uno dei termini noti $b'_i$ con $i > r$ è diverso da $0$ (mentre la riga a sinistra è di zeri), il sistema è IMPOSSIBILE.
    \[ 0 = b'_i \neq 0 \implies \text{Nessuna soluzione} \]
    
    \item \textbf{Caso Risolubile:} Se $b'_i = 0$ per ogni $i > r$, allora le ultime $m-r$ righe sono identità ($0=0$) e non danno informazioni (possono essere trascurate).
    \begin{itemize}
        \item Il sistema ammette soluzioni.
        \item Numero di variabili libere = $n - r$.
        \item Soluzioni totali: $\infty^{n-r}$.
    \end{itemize}
\end{enumerate}
\end{concept}

% -------------------------------------------------------------------------
\section{Esempio Pratico e Raffinamento della Matrice}

Consideriamo un esempio di matrice associata a un sistema, dove $C$ è univocamente determinata:

\[
\begin{pmatrix}
1 & 5 & 2 & 0 & 7 & | & 1 \\
0 & 0 & 4 & 1 & 2 & | & 1 \\
0 & 0 & 0 & 2 & 0 & | & 0 \\
0 & 0 & 0 & 0 & 1 & | & 7
\end{pmatrix}
\]
I pivot sono cerchiati (qui indicati dalle posizioni: $1, 4, 2, 1$).
Le colonne senza pivot corrispondono alle \textbf{Variabili Libere}.

\subsection{Strategie di riduzione (Domande dalla lavagna)}

\begin{observation}{Fase 1: Normalizzazione dei Pivot}{}
\textit{"Posso trasformare la matrice ridotta in una matrice equivalente in cui $P_i = 1, \forall i$?"} \\
\textbf{Sì.} È sufficiente dividere ogni riga per il proprio pivot.
Esempio di trasformazione sulla matrice sopra:
\[
\begin{pmatrix}
1 & 5 & 2 & 0 & 7 & | & 1 \\
0 & 0 & 1 & 1/4 & 1/2 & | & 1/4 \\
0 & 0 & 0 & 1 & 0 & | & 5/2 \\
0 & 0 & 0 & 0 & 1 & | & 7
\end{pmatrix}
\]
\end{observation}

\begin{observation}{Fase 2: Gauss-Jordan (Zeri sopra i Pivot)}{}
\textit{"Posso trasformare la matrice ottenuta in (1) in una matrice che ha tutti zeri non solo sotto i pivot, ma anche sopra?"} \\
\textbf{Sì.} Si opera "all'indietro" (dal basso verso l'alto).
Esempio di operazione mostrato alla lavagna:
\[
C_1 \to C_1 - 5C_2 \quad (\text{Notazione usata, probabile riferimento a operazioni su righe/colonne per eliminazione})
\]
Obiettivo: Ottenere una forma diagonale (o simile) per le colonne dei pivot.
\end{observation}

\restoregeometry

\newpage
\newgeometry{top=2.5cm, bottom=2.5cm, left=2.5cm, right=2.5cm}

\section{Approfondimento: Assiomi e Spazio dei Polinomi}

\subsection{Definizione Generale di Spazio Vettoriale (Dalle slide)}

\begin{definition}{Spazio Vettoriale su $\mathbb{R}$}{}
Uno \textbf{Spazio Vettoriale su $\mathbb{R}$} è un insieme $V$ dotato di due operazioni:
\begin{itemize}
    \item \textbf{Somma:} $+ : V \times V \to V$
    \item \textbf{Prodotto per scalare:} $\cdot : \mathbb{R} \times V \to V$
\end{itemize}
Tali che valgano le seguenti 8 proprietà (assiomi).
\end{definition}

\begin{enumerate}
    \item $(u + v) + w = u + (v + w)$ [Associativa della somma]
    \item $\exists \, 0 \in V$ tale che $0 + v = v + 0 = v, \quad \forall v \in V$ [Esistenza dell'elemento neutro]
    \item $\forall v \in V, \exists -v \in V : v + (-v) = 0$ [Esistenza dell'opposto]
    \item $\forall v, w \in V \quad v + w = w + v$ [Commutativa della somma]
    \item $(\lambda \mu) v = \lambda (\mu v)$ [Associativa del prodotto per scalare]
    \item $1 \cdot v = v$ e $0 \cdot v = 0 \quad \forall v \in V$ [Elemento neutro moltiplicativo / annullamento]
    \item $\lambda (v + w) = \lambda v + \lambda w \quad \forall \lambda \in \mathbb{R}, \forall v, w \in V$ [Distributiva rispetto alla somma di vettori]
    \item $(\lambda + \mu) v = \lambda v + \mu v \quad \forall \lambda, \mu \in \mathbb{R}, \forall v \in V$ [Distributiva rispetto alla somma di scalari]
\end{enumerate}

\section{Esempio 1: Spazio dei Polinomi (Dalla lavagna)}

\begin{definition}{Polinomi $P_n(x)$}{}
I Polinomi di grado $n$ in una variabile $x$ si indicano con $P_n(x)$.
\[
P_n(x) = \{ a_n x^n + a_{n-1} x^{n-1} + \dots + a_1 x + a_0 \mid a_i \in \mathbb{R} \quad \forall i \}
\]
\end{definition}

Per capire se posso dotarlo di una struttura di spazio vettoriale, devo innanzitutto definire le due operazioni e poi verificare che soddisfino le proprietà (la lista degli assiomi vista sopra).

\begin{proposition}{Definizione delle Operazioni}{}
Siano $p(x), q(x) \in P_n(x)$ dove:
\[
\begin{aligned}
p(x) &= a_n x^n + \dots + a_0 \\
q(x) &= b_n x^n + \dots + b_0
\end{aligned}
\]

\begin{itemize}
    \item \textbf{Somma ($+$):} $P_n(x) \times P_n(x) \to P_n(x)$ \\
    $(p(x), q(x)) \to p(x) + q(x) = (a_n + b_n)x^n + \dots + (a_0 + b_0)$ \\
    (Si sommano i coefficienti corrispondenti)
    
    \item \textbf{Prodotto per scalare ($\cdot$):} $\mathbb{R} \times P_n(x) \to P_n(x)$ \\
    $(\lambda, p(x)) \to \lambda \cdot p(x) = (\lambda a_n)x^n + \dots + (\lambda a_0)$
\end{itemize}
\end{proposition}

\subparagraph{Dimostrazione alla lavagna (Verifica degli assiomi):}
Il professore dimostra che $(P_n(x), +, \cdot)$ è uno spazio vettoriale su $\mathbb{R}$.

\textbf{I. Associatività della somma:}\\
Siano $p(x), q(x), h(x) \in P_n(x)$. Bisogna dimostrare che $(p(x)+q(x))+h(x) = p(x)+(q(x)+h(x))$.

Svolgimento:
\[
\begin{aligned}
&[(a_n+b_n)x^n + \dots + (a_0+b_0)] + h(x) \\
&= [(a_n+b_n) + c_n]x^n + \dots + [(a_0+b_0) + c_0]
\end{aligned}
\]
(Nota scritta a mano: "$||$ perché siamo in $\mathbb{R}$" - intende che l'associatività vale per i numeri reali)
\[
= [a_n + (b_n+c_n)]x^n + \dots + [a_0 + (b_0+c_0)]
\]
Che corrisponde a $p(x) + (q(x) + h(x))$.

\textbf{II. Elemento Neutro:}\\
$\exists \, 0 \in P_n(x)$ tale che $0 + p(x) = p(x) \quad \forall p \in P_n(x)$.\\
L'elemento è il polinomio con tutti i coefficienti nulli: $0 = 0x^n + \dots + 0$.

\begin{note}{Note sulle altre proprietà (Lavagna)}{}
\begin{itemize}
    \item $p(x) \in P_n(x) \implies \exists -p(x) \in P_n(x)$ tale che $p(x) + (-p(x)) = 0$. (Esistenza dell'opposto, coefficienti $-a_i$).
    \item Commutativa: $p(x) + q(x) = q(x) + p(x)$ (deriva dalla commutatività di $a_i + b_i$ in $\mathbb{R}$).
    \item $(\lambda \mu)p(x) = \lambda (\mu p(x))$
    \item Elemento neutro per il prodotto scalare: $1 \in \mathbb{R}$.
    \item Nota finale: "7 e 8 si dimostrano analogamente".
\end{itemize}
\end{note}

\section{Esempio 2: Vettori Geometrici (Dalla slide)}

\textbf{Contesto:} Consideriamo il piano euclideo $\mathbb{R}^2$ e fissiamo un'origine $O \in \mathbb{R}^2$.

\begin{definition}{Insieme $V_O^2$}{}
Un \textbf{Vettore Applicato in $O$} $\doteq$ un segmento orientato con primo estremo in $O$ e secondo estremo in un punto $A \in \mathbb{R}^2$.\\
Si indica con $\overrightarrow{OA}$.

$V_O^2 \doteq \{ \text{vettori applicati in } O \}$\\
L'Origine di $V_O^2$ è il punto $O$ (vettore nullo).
\end{definition}

\begin{proposition}{Operazioni Geometriche}{}
\begin{itemize}
    \item \textbf{Somma di vettori applicati ($+$):} $V_O^2 \times V_O^2 \to V_O^2$ \\
    Dati $\overrightarrow{OA}$ e $\overrightarrow{OB}$, la somma è $\overrightarrow{OC}$ tale che $\overrightarrow{OA} + \overrightarrow{OB} = \overrightarrow{OC}$.\\
    (\textbf{Regola del Parallelogramma}: $C$ è il vertice opposto a $O$ del parallelogramma formato da $A$ e $B$).

    \item \textbf{Prodotto per scalare ($\cdot$):} $\mathbb{R} \times V_O^2 \to V_O^2$ \\
    $(\lambda, \overrightarrow{OA}) \to \lambda \cdot \overrightarrow{OA} = \text{un vettore } \overrightarrow{OC}$ dove:
    \begin{itemize}
        \item $C$ appartiene alla retta per $O$ ed $A$.
        \item Tale che $\frac{\text{lunghezza } OC}{\text{lunghezza } OA} = |\lambda|$.
    \end{itemize}
\end{itemize}
\end{proposition}

\begin{tcolorbox}[colback=white, colframe=cherryRed, title=Affermazione (AFF)]
$(V_O^2, +, \text{prod per scalare})$ è uno spazio vettoriale.
\end{tcolorbox}

\restoregeometry

\newpage
\section{Spazi Vettoriali}

\begin{concept}{Definizione di Spazio Vettoriale}
Uno spazio vettoriale $V$ su un campo $K$ è un insieme dotato di due operazioni (somma e prodotto per scalare) che soddisfano 8 assiomi.
\end{concept}

\subsection{Basi e Dimensione}
\begin{itemize}
    \item Generatori
    \item Indipendenza Lineare
    \item Teorema della Base
\end{itemize}

\newpage
\newgeometry{top=2.5cm, bottom=2.5cm, left=2cm, right=2cm}
\setlength{\parindent}{0pt}
\setlength{\parskip}{6pt}

% --- TITOLO ---
\begin{center}
    {\Huge \textbf{\color{cherryRed} Vettori Applicati nel Piano}} \\
    \vspace{0.3cm}
    {\Large \textit{Struttura dello Spazio Vettoriale $V_O^2$}}
    \vspace{0.5cm}
    \hrule height 1pt
\end{center}

\section{Costruzione Geometrica di $V_O^2$}

Dalla proiezione in aula (slide 3), costruiamo lo spazio vettoriale basato sulla geometria euclidea.

\begin{definition}{Vettore Applicato in $O$}{}
Consideriamo il piano euclideo $\mathbb{R}^2$ e fissiamo un punto $O \in \mathbb{R}^2$ detto \textbf{Origine} (o punto fissato).

Un \textbf{Vettore Applicato in $O$} è definito come un segmento orientato con:
\begin{itemize}
    \item \textbf{Primo estremo:} il punto $O$.
    \item \textbf{Secondo estremo:} un punto $A \in \mathbb{R}^2$.
\end{itemize}
Si denota con $\overrightarrow{OA}$.
\end{definition}

\subsection{L'Insieme $V_O^2$}
Definiamo l'insieme di tutti i vettori applicati nell'origine:
\[
V_O^2 \doteq \{ \text{vettori applicati in } O \}
\]
L'\textbf{Origine} di $V_O^2$ corrisponde al vettore nullo, ovvero il punto $O$ stesso (segmento degenere).

% =========================================================================
\section{Operazioni in $V_O^2$}

Per dotare l'insieme $V_O^2$ di una struttura di spazio vettoriale, definiamo le due operazioni fondamentali.

\subsection{1. Somma di Vettori}
L'operazione somma è definita come:
\[ + : V_O^2 \times V_O^2 \longrightarrow V_O^2 \]
Dati due vettori $\overrightarrow{OA}$ e $\overrightarrow{OB}$, la somma è il vettore $\overrightarrow{OC}$ tale che $C$ è il quarto vertice del parallelogramma costruito su $OA$ e $OB$.

\[ \overrightarrow{OA} + \overrightarrow{OB} = \overrightarrow{OC} \]

\begin{center}
\begin{tikzpicture}[>=Stealth, scale=1.5]
    % Coordinate
    \coordinate (O) at (0,0);
    \coordinate (A) at (2, 0.5);
    \coordinate (B) at (1, 1.5);
    \coordinate (C) at (3, 2); % A + B

    % Vettori
    \draw[->, thick, color=maroon] (O) -- (A) node[midway, below right] {$\overrightarrow{OA}$};
    \draw[->, thick, color=maroon] (O) -- (B) node[midway, left] {$\overrightarrow{OB}$};
    \draw[->, ultra thick, color=cherryRed] (O) -- (C) node[midway, above left] {$\overrightarrow{OC} = \text{Somma}$};

    % Linee tratteggiate per parallelogramma
    \draw[dashed, gray] (A) -- (C);
    \draw[dashed, gray] (B) -- (C);

    % Etichette Punti
    \node[below left] at (O) {$O$};
    \node[right] at (A) {$A$};
    \node[above] at (B) {$B$};
    \node[above right] at (C) {$C$};
\end{tikzpicture}
\end{center}

\subsection{2. Prodotto per Scalare}
L'operazione prodotto per uno scalare $\lambda \in \mathbb{R}$ è definita come:
\[ \cdot : \mathbb{R} \times V_O^2 \longrightarrow V_O^2 \]
Dato uno scalare $\lambda$ e un vettore $\overrightarrow{OA}$, il risultato è un vettore $\overrightarrow{OC}$ tale che:
\begin{itemize}
    \item $C$ appartiene alla retta passante per $O$ e per $A$.
    \item La lunghezza (modulo) soddisfa:
    \[ \frac{\text{lunghezza } OC}{\text{lunghezza } OA} = |\lambda| \]
\end{itemize}
Se $\lambda > 0$, il verso è lo stesso. Se $\lambda < 0$, il verso è opposto.

\textbf{Esempio dalla lavagna ($\lambda = 2$):}

\begin{center}
\begin{tikzpicture}[>=Stealth, scale=1.5]
    % Coordinate
    \coordinate (O) at (0,0);
    \coordinate (A) at (1.5, 0.5);
    \coordinate (C) at (3, 1); % 2 * A

    % Retta tratteggiata
    \draw[dashed, gray] (-0.5, -0.16) -- (3.5, 1.16);

    % Vettori
    \draw[->, thick, color=maroon] (O) -- (A) node[midway, above left] {$\overrightarrow{OA}$};
    \draw[->, ultra thick, color=cherryRed, transform canvas={yshift=-2pt}] (O) -- (C) node[midway, below right] {$\lambda \cdot \overrightarrow{OA}$};

    % Etichette
    \node[below left] at (O) {$O$};
    \node[above] at (A) {$A$};
    \node[above] at (C) {$C$};
    \node[right] at (3.5, 1) {$\lambda=2$};
\end{tikzpicture}
\end{center}

\begin{tcolorbox}[colback=white, colframe=cherryRed, title=Affermazione (AFF)]
L'insieme $V_O^2$, dotato delle operazioni di somma (regola del parallelogramma) e prodotto per scalare (dilatazione/contrazione), è uno \textbf{Spazio Vettoriale} su $\mathbb{R}$.
\end{tcolorbox}

\begin{observation}{Scopo della Lezione (Blackboard)}{}
Alla base della lavagna si legge:
\textit{"Scopo: Dotare $V_O^2$ di una struttura di spazio vettoriale."} 
Questo significa verificare che le operazioni geometriche appena definite soddisfino gli 8 assiomi visti precedentemente (associatività, elemento neutro, ecc.).
\end{observation}

\restoregeometry

\newpage
\section{Applicazioni Lineari}

\begin{definitionbox}
Una funzione $f: V \to W$ tra due spazi vettoriali è \textbf{lineare} se conserva le operazioni, ovvero se $\forall v, w \in V, \forall \lambda \in \mathbb{R}$:
1) $f(v+w) = f(v) + f(w)$
2) $f(\lambda v) = \lambda f(v)$
\end{definitionbox}

\begin{example}{Esempi di Applicazioni}{}
\textbf{1.} $f(x) = 3x$ \textbf{è lineare}. (Vale per $f(x) = kx \ \forall k$).\\
\textbf{2.} $f(x) = ax + b$ \textbf{NON è lineare} se $b \neq 0$.\\
\textbf{3.} $f(x) = x^2$ \textbf{NON è lineare}.
\end{example}

\begin{concept}{Applicazione Associata a una Matrice}
Consideriamo $A \in M_{m \times n}(\mathbb{R})$ e l'applicazione $L_A: \mathbb{R}^n \to \mathbb{R}^m$ definita da:
$$ L_A(x) = A \cdot x $$
$L_A$ è lineare. L'immagine di $e_i$ è l'$i$-esima colonna di $A$: $L_A(e_i) = A^i$.
Tutte le applicazioni lineari tra spazi reali finitamente generati sono matrici in incognito.
\end{concept}

\subsection{Costruzione di un'Applicazione Lineare}
\begin{proposition}{Esistenza ed Unicità}{}
Per costruire un'applicazione lineare $L: V \to W$ è sufficiente fissare i valori su una QUALSIASI BASE fissata per $V$. Siano $v_1, \dots, v_n$ base di $V$ e $w_1, \dots, w_n \in W$ a piacere, esiste un'unica applicazione lineare tale che $L(v_i) = w_i$.
\end{proposition}

\section{Nucleo, Immagine e Teorema della Dimensione}

\begin{definitionbox}
Sia $f: V \to W$ un'applicazione lineare. 
\begin{itemize}
    \item \textbf{Nucleo:} $Ker(f) = \{ v \in V \mid f(v) = 0_W \}$. È un sottospazio vettoriale di $V$.
    \item \textbf{Immagine:} $Im(f) = \{ w \in W \mid \exists v \in V : f(v) = w \}$. È un sottospazio vettoriale di $W$.
\end{itemize}
\end{definitionbox}

\begin{proposition}{Criterio di Iniettività}{}
$f$ è iniettiva $\iff Ker(f) = \{ 0 \}$. (i.e. il nucleo contiene solo il vettore nullo).
\end{proposition}

\begin{concept}{Nucleo e Immagine per Matrici}
Chi è il $Ker(L_A)$?\\
È l'insieme delle soluzioni del sistema omogeneo associato alla matrice $A$. $L_A$ è iniettiva $\iff$ le colonne di $A$ sono linearmente indipendenti.\\
Chi è l'$Im(L_A)$?\\
L'immagine di $L_A$ è lo span delle colonne di $A$. Di conseguenza, la dimensione dell'immagine è il rango della matrice: $\dim Im(L_A) = \text{rk}(A)$.
\end{concept}

\begin{proposition}{Teorema della Dimensione (Rank-Nullity)}{}
Sia $f: V \to W$ un'applicazione lineare. Se $V$ ha dimensione finita $n$, allora:
$$ \dim Ker(f) + \dim Im(f) = \dim(V) = n $$
\end{proposition}

\newpage
\section{Autovalori e Diagonalizzazione}

\begin{concept}{Autovalore e Autovettore}
Un vettore $v \neq 0$ è autovettore di $f$ relativo all'autovalore $\lambda$ se $f(v) = \lambda v$.
\end{concept}

\subsection{Polinomio Caratteristico}
$P_A(\lambda) = \det(A - \lambda I)$

\newpage
\section{Prodotti Scalari e Ortogonalità}

\begin{concept}{Prodotto Scalare Standard}
In $\mathbb{R}^n$, $\langle x, y \rangle = \sum_{i=1}^n x_i y_i$.
\end{concept}

\subsection{Procedimento di Gram-Schmidt}
Metodo per ortonormalizzare una base.

\newpage
\newpage
\newgeometry{top=1.5cm, bottom=1.5cm, left=1cm, right=1cm} % Margini ridotti per Cheat Sheet
\section*{Cheat Sheet Finale: Algebra Lineare}
\addcontentsline{toc}{section}{Cheat Sheet Finale}

\begin{multicols*}{2}

% --- SISTEMI LINEARI ---
\begin{tcolorbox}[colback=lightBg, colframe=cherryRed, title={Sistemi Lineari ($Ax=b$)}]
\small
\textbf{Rouché-Capelli:}
\begin{itemize}[leftmargin=*]
    \item Compatibile $\iff \text{rk}(A) = \text{rk}(A|b)$.
    \item $\infty^{n-r}$ soluzioni se $r < n$.
    \item 1 soluzione unica se $r = n$.
\end{itemize}
\textbf{Gauss:}
\begin{enumerate}[leftmargin=*]
    \item Scambio righe ($R_i \leftrightarrow R_j$).
    \item Moltiplicazione per scalare ($\lambda R_i$).
    \item Comb. lineare ($R_i \to R_i + k R_j$).
\end{enumerate}
\textbf{Cramer:} Se $\det(A) \neq 0, x_i = \frac{\det(A_i)}{\det(A)}$.
\end{tcolorbox}

% --- SPAZI VETTORIALI ---
\begin{tcolorbox}[colback=white, colframe=maroon, title={Spazi Vettoriali e Basi}]
\small
\textbf{Indipendenza Lineare:} $\sum \lambda_i v_i = 0 \implies \forall \lambda_i = 0$.
\textbf{Generatori:} $\text{Span}(v_1, \dots, v_k) = V$.
\textbf{Base:} Insieme di generatori lin. indip.
\textbf{Dimensione:} Numero elementi della base.
\tcblower
\textbf{Formula di Grassmann:}
\[ \dim(U+W) = \dim(U) + \dim(W) - \dim(U \cap W) \]
\textbf{Somma Diretta ($V = U \oplus W$):}
\begin{itemize}[leftmargin=*]
    \item $U \cap W = \{0\}$.
    \item Ogni $v \in V$ si scrive in modo unico come $u+w$.
\end{itemize}
\end{tcolorbox}

% --- MATRICI ---
\begin{tcolorbox}[colback=lightBg, colframe=cherryRed, title={Matrici e Determinante}]
\small
\textbf{Rango (rk):} Numero pivot / Numero righe lin. indip.
\textbf{Inversa ($A^{-1}$):} Esiste $\iff \det(A) \neq 0$.
\[ A^{-1} = \frac{1}{\det(A)} (\text{Cof}(A))^T \]
\textbf{Proprietà Determinante:}
\begin{itemize}[leftmargin=*]
    \item $\det(A) = \det(A^T)$.
    \item $\det(AB) = \det(A)\det(B)$ (Binet).
    \item $\det(A^{-1}) = 1/\det(A)$.
\end{itemize}
\end{tcolorbox}

% --- APPLICAZIONI LINEARI ---
\begin{tcolorbox}[colback=white, colframe=maroon, title={Applicazioni Lineari ($f: V \to W$)}]
\small
\textbf{Nucleo ($\ker f$):} $\{ v \in V \mid f(v) = 0 \}$. (Iniettiva $\iff \ker f = \{0\}$).
\textbf{Immagine ($\text{Im } f$):} $\{ w \in W \mid \exists v: f(v)=w \}$. (Suriettiva $\iff \text{Im } f = W$).
\tcblower
\textbf{Teorema della Dimensione (Rank-Nullity):}
\[ \dim(V) = \dim(\ker f) + \dim(\text{Im } f) \]
\textbf{Matrice Associata ($M_\mathcal{B}^\mathcal{C}(f)$):}
Le colonne sono le coordinate delle immagini dei vettori della base di partenza rispetto alla base di arrivo.
\[ [f(v)]_\mathcal{C} = A [v]_\mathcal{B} \]
\end{tcolorbox}

% --- AUTOVALORI ---
\begin{tcolorbox}[colback=lightBg, colframe=cherryRed, title={Autovalori e Diagonalizzazione}]
\small
\textbf{Definizione:} $f(v) = \lambda v \implies (A - \lambda I)v = 0$.
\textbf{Polinomio Caratteristico:} $P_A(\lambda) = \det(A - \lambda I)$.
\textbf{Spettro:} Insieme delle radici di $P_A(\lambda)$.
\textbf{Molteplicità:}
\begin{itemize}[leftmargin=*]
    \item Alg. ($m_a(\lambda)$): esponente nel pol. caratt.
    \item Geom. ($m_g(\lambda)$): $\dim(\ker(A-\lambda I)) = n - \text{rk}(A-\lambda I)$.
\end{itemize}
\textbf{Diagonalizzabile} $\iff$:
\begin{itemize}[leftmargin=*]
    \item Tutte le radici sono nel campo (p.c. si spezza).
    \item $\forall \lambda, m_g(\lambda) = m_a(\lambda)$.
\end{itemize}
Se diag., $D = P^{-1}AP$ con $P$ matrice cambio base (autovettori).
\end{tcolorbox}

% --- PRODOTTI SCALARI ---
\begin{tcolorbox}[colback=white, colframe=maroon, title={Prodotti Scalari e Spettale}]
\small
\textbf{Standard in $\mathbb{R}^n$:} $\langle x, y \rangle = x^T y$.
\textbf{Norma:} $\|v\| = \sqrt{\langle v, v \rangle}$.
\textbf{Ortogonalità:} $\langle v, w \rangle = 0$.
\textbf{Proiezione Ortogonale su $v$:} $P_v(u) = \frac{\langle u, v \rangle}{\|v\|^2} v$.
\textbf{Gram-Schmidt:} $u_{k} = v_k - \sum_{j=1}^{k-1} P_{u_j}(v_k)$.
\tcblower
\textbf{Teorema Spettrale (Reale):}
Una matrice $A \in \mathbb{R}^{n \times n}$ è simmetrica ($A=A^T$) $\iff$ ammette una base ORTONORMALE di autovettori.
(Tutte le matrici simmetriche sono diagonalizzabili).
\end{tcolorbox}

\end{multicols*}
\restoregeometry


\end{document}